% Implementação do abastecimento de água, do começo ao fim
% (leitura em zigue-zague: linhas ímpares da esquerda para a direita,
%  linhas pares da direita para a esquerda)
\begin{tikzpicture}[
    passo/.style={etapa=2.8cm, minimum height=1.1cm},
    nomefaixa/.style={font=\scriptsize\bfseries, text=azulpsh, align=left,
                      text width=3.4cm, anchor=west}
  ]
  % 4 colunas (x = 5.3, 8.7, 12.1, 15.5) e 6 linhas espaçadas de 1.8 cm
  % Linha 1 -- esquerda para a direita
  \node[passo] (a1) at (5.3,0)  {Planejamento integrado das ações};
  \node[passo] (a2) at (8.7,0)  {Seleção dos municípios pelos critérios};
  \node[passo] (a3) at (12.1,0)  {Articulação e adesão dos municípios};
  \node[passo] (a4) at (15.5,0) {Termo de Cooperação com o município};
  % Linha 2 -- direita para a esquerda
  \node[passo] (b1) at (15.5,-1.8) {Modelo de gestão e pré-adesão da comunidade};
  \node[passo] (b2) at (12.1,-1.8)  {Mobilização social e busca ativa};
  \node[passo] (b3) at (8.7,-1.8)  {Diagnóstico e enquadramento das famílias};
  \node[passo] (b4) at (5.3,-1.8)  {Termo de adesão das famílias};
  % Linha 3 -- esquerda para a direita
  \node[passo] (c1) at (5.3,-3.6)  {Estudo hidrogeológico, locação e anuência};
  \node[passo] (c2) at (8.7,-3.6)  {Perfuração do poço};
  \node[passo] (c3) at (12.1,-3.6)  {Teste de vazão e análise físico-química};
  \node[passo] (c4) at (15.5,-3.6) {Outorga de direito de uso da água};
  % Linha 4 -- direita para a esquerda
  \node[passo] (d1) at (15.5,-5.4) {Anteprojeto da solução};
  \node[passo] (d2) at (12.1,-5.4)  {Projeto básico e aprovação};
  \node[passo] (d3) at (8.7,-5.4)  {Projeto executivo};
  \node[passo] (d4) at (5.3,-5.4)  {Estudos e licenças ambientais};
  % Linha 5 -- esquerda para a direita
  \node[passo] (e1) at (5.3,-7.2)  {Obras: casa química, reservatório e distribuição};
  \node[passo] (e2) at (8.7,-7.2)  {Supervisão e fiscalização das obras};
  \node[passo] (e3) at (12.1,-7.2)  {\textit{As built}, pré-operação e testes};
  \node[passo] (e4) at (15.5,-7.2) {Capacitação dos operadores e usuários};
  % Linha 6 -- direita para a esquerda
  \node[passo] (f1) at (15.5,-9.0) {Reunião técnica de entrega};
  \node[passo] (f2) at (12.1,-9.0)  {Aceite da entidade gestora};
  \node[passo] (f3) at (8.7,-9.0)  {Operação e suporte pós-implantação};
  \node[marco=2.8cm, minimum height=1.1cm] (f4) at (5.3,-9.0) {TRD e encerramento};

  % faixas e nomes das fases
  \begin{scope}[on background layer]
    \foreach \y in {0,-1.8,-3.6,-5.4,-7.2,-9.0}
      \fill[cinzaclaro!70, rounded corners=2pt] (-0.2,\y+0.8) rectangle (17.4,\y-0.8);
  \end{scope}
  \node[nomefaixa] at (-0.1,0)    {1. Preparação\\ institucional};
  \node[nomefaixa] at (-0.1,-1.8) {2. Gestão e\\ mobilização};
  \node[nomefaixa] at (-0.1,-3.6) {3. Captação};
  \node[nomefaixa] at (-0.1,-5.4) {4. Projeto e\\ licenciamento};
  \node[nomefaixa] at (-0.1,-7.2) {5. Obras};
  \node[nomefaixa] at (-0.1,-9.0) {6. Entrega e\\ operação};

  % ligações
  \foreach \l in {a,c,e} {
    \draw[seta] (\l1) -- (\l2); \draw[seta] (\l2) -- (\l3); \draw[seta] (\l3) -- (\l4);
  }
  \foreach \l in {b,d,f} {
    \draw[seta] (\l1) -- (\l2); \draw[seta] (\l2) -- (\l3); \draw[seta] (\l3) -- (\l4);
  }
  \draw[seta] (a4) -- (b1);
  \draw[seta] (b4) -- (c1);
  \draw[seta] (c4) -- (d1);
  \draw[seta] (d4) -- (e1);
  \draw[seta] (e4) -- (f1);
\end{tikzpicture}
